% Focused custom theorem style/counter matrix.
% Build from thesis/ after the default formats have initialized, then create
% unique environments so the public customization contract can be exercised.
\input{./template/configure}

\newcounter{MatrixCustomCounter}

% Freeze dynamic registry parser expansion, per-call defaults, and omission.
\def\NCKUTheoremEnv{EnvExpandedDefinition}
\def\NCKUTheoremText{Expanded Definition}
\def\NCKUTheoremCounter{MatrixCustomCounter}
\NCKUPrivateSetTheoremFormatKeys{Definition}{EnvironmentName=\NCKUTheoremEnv,ShowText=\NCKUTheoremText,FollowCounter=\NCKUTheoremCounter}
\def\NCKUTheoremEnv{CHANGED}
\def\NCKUTheoremText{CHANGED}
\def\NCKUTheoremCounter{CHANGED}
\typeout{NCKU-THEOREM-REGISTRY-EXPANDED: \GetTheoremDefinitionFormatEnvironmentName/\GetTheoremDefinitionFormatShowText/\GetTheoremDefinitionFormatFollowCounter}
\NCKUPrivateSetTheoremFormatKeys{Definition}{ShowText=Partial Definition}
\typeout{NCKU-THEOREM-REGISTRY-PARTIAL: \GetTheoremDefinitionFormatEnvironmentName/\GetTheoremDefinitionFormatShowText/\GetTheoremDefinitionFormatFollowCounter}
\NCKUPrivateSetTheoremFormatKeys{Definition}{\empty}
\typeout{NCKU-THEOREM-REGISTRY-OMITTED: \GetTheoremDefinitionFormatEnvironmentName/\GetTheoremDefinitionFormatShowText/\GetTheoremDefinitionFormatFollowCounter}
\NCKUPrivateSetTheoremFormatKeys{Note}{\empty}
\typeout{NCKU-THEOREM-REGISTRY-OPTIONAL: \GetTheoremNoteFormatEnvironmentName/\GetTheoremNoteFormatShowText/\GetTheoremNoteFormatFollowCounter}

\newcommand{\ConfigureMatrixTheorem}[2]{%
  \SetTheoremFormat[#1]{%
    EnvironmentName = {EnvMatrix#1},
    ShowText = {Matrix #1},
    FollowCounter = #2,
  }%
}

% Configure every type before aggregate initialization. This deliberately makes
% Problem/Example forward-reference Theorem, whose own two-hop chain resolves
% through Definition to an empty effective parent.
\ConfigureMatrixTheorem{Definition}{\empty}
\ConfigureMatrixTheorem{Theorem}{Definition}
\ConfigureMatrixTheorem{Condition}{Section}
\ConfigureMatrixTheorem{Problem}{Theorem}
\ConfigureMatrixTheorem{Example}{Theorem}
\ConfigureMatrixTheorem{Lemma}{Theorem}
\ConfigureMatrixTheorem{Corollary}{Theorem}
\ConfigureMatrixTheorem{Proposition}{Theorem}
\ConfigureMatrixTheorem{Conjecture}{Theorem}
\ConfigureMatrixTheorem{Criterion}{Theorem}
\ConfigureMatrixTheorem{Assertion}{Theorem}
% Preserve arbitrary existing LaTeX counters as terminal values, including a
% registered theorem chain that resolves to the custom counter.
\ConfigureMatrixTheorem{Question}{MatrixCustomCounter}
\ConfigureMatrixTheorem{Hypothesis}{Question}

% Optional-numbering forms keep starred behavior when empty, but Note follows
% the unscoped theorem chain and Summary explicitly follows Section.
\ConfigureMatrixTheorem{Proof}{\empty}
\ConfigureMatrixTheorem{Note}{Theorem}
\ConfigureMatrixTheorem{Annotation}{\empty}
\ConfigureMatrixTheorem{Claim}{\empty}
\ConfigureMatrixTheorem{Case}{\empty}
\ConfigureMatrixTheorem{Acknowledgment}{\empty}
\ConfigureMatrixTheorem{Conclusion}{\empty}
\ConfigureMatrixTheorem{Summary}{Section}

\InitTheoremFormats
\setcounter{MatrixCustomCounter}{3}

\newcommand{\TypeoutMatrixTheorem}[1]{%
  \typeout{NCKU-TEST-MATRIX-#1: \csname GetTheorem#1FormatEnvironmentName\endcsname/\csname GetTheorem#1FormatShowText\endcsname/\csname GetTheorem#1FormatFollowCounter\endcsname}%
}

\begin{document}
\TypeoutMatrixTheorem{Definition}
\TypeoutMatrixTheorem{Theorem}
\TypeoutMatrixTheorem{Condition}
\TypeoutMatrixTheorem{Problem}
\TypeoutMatrixTheorem{Example}
\TypeoutMatrixTheorem{Lemma}
\TypeoutMatrixTheorem{Corollary}
\TypeoutMatrixTheorem{Proposition}
\TypeoutMatrixTheorem{Conjecture}
\TypeoutMatrixTheorem{Criterion}
\TypeoutMatrixTheorem{Assertion}
\TypeoutMatrixTheorem{Question}
\TypeoutMatrixTheorem{Hypothesis}
\TypeoutMatrixTheorem{Proof}
\TypeoutMatrixTheorem{Note}
\TypeoutMatrixTheorem{Annotation}
\TypeoutMatrixTheorem{Claim}
\TypeoutMatrixTheorem{Case}
\TypeoutMatrixTheorem{Acknowledgment}
\TypeoutMatrixTheorem{Conclusion}
\TypeoutMatrixTheorem{Summary}

\stepcounter{section}
\InsertDefinition{MATRIXDEFINITIONBODYONE}
\InsertTheorem{MATRIXTHEOREMBODYONE}
\InsertCondition{MATRIXCONDITIONBODYONE}
\InsertProblem{MATRIXPROBLEMBODYONE}
\InsertExample{MATRIXEXAMPLEBODYONE}
\InsertLemma{MATRIXLEMMABODYONE}
\InsertCorollary{MATRIXCOROLLARYBODYONE}
\InsertProposition{MATRIXPROPOSITIONBODYONE}
\InsertConjecture{MATRIXCONJECTUREBODYONE}
\InsertCriterion{MATRIXCRITERIONBODYONE}
\InsertAssertion{MATRIXASSERTIONBODYONE}
\InsertQuestion{MATRIXQUESTIONBODYONE}
\InsertHypothesis{MATRIXHYPOTHESISBODYONE}
\InsertProof{MATRIXPROOFBODYONE}
\InsertNote{MATRIXNOTEBODYONE}
\InsertAnnotation{MATRIXANNOTATIONBODYONE}
\InsertClaim{MATRIXCLAIMBODYONE}
\InsertCase{MATRIXCASEBODYONE}
\InsertAcknowledgment{MATRIXACKNOWLEDGMENTBODYONE}
\InsertConclusion{MATRIXCONCLUSIONBODYONE}
\InsertSummary{MATRIXSUMMARYBODYONE}

\stepcounter{section}
\InsertDefinition{MATRIXDEFINITIONBODYTWO}
\InsertTheorem{MATRIXTHEOREMBODYTWO}
\InsertCondition{MATRIXCONDITIONBODYTWO}
\InsertNote{MATRIXNOTEBODYTWO}
\InsertSummary{MATRIXSUMMARYBODYTWO}

\typeout{NCKU-TEST-PASS: custom theorem style/counter matrix compiled}
\end{document}
